\documentclass[12pt]{article}
\usepackage[margin=1in, top=0.9in, bottom=0.9in]{geometry}
\usepackage{newtxtext}
\usepackage{xcolor}
\usepackage{enumitem}
\usepackage{booktabs}
\usepackage{microtype}
\usepackage[colorlinks=false, hidelinks]{hyperref}
\usepackage{fancyhdr}
\usepackage{mdframed}

% Colors
\definecolor{pennblue}{HTML}{011F5B}
\definecolor{pennred}{HTML}{990000}
\definecolor{lightgray}{HTML}{F5F5F5}
\definecolor{slidegray}{HTML}{DDDDDD}

% Header / Footer
\pagestyle{fancy}
\fancyhf{}
\renewcommand{\headrulewidth}{0.4pt}
\fancyhead[L]{\small\color{pennblue}\textbf{Cross-Dataset HAR Transfer} — Speaking Script}
\fancyhead[R]{\small\color{pennblue}Gengzhan · Forest · Gengzhan}
\fancyfoot[C]{\small\thepage}

% Section styling — slide labels
\newcommand{\slidetitle}[1]{%
  \vspace{14pt}%
  \noindent\colorbox{pennblue}{\parbox{\dimexpr\textwidth-2\fboxsep\relax}{%
    \strut\bfseries\color{white}\quad #1\strut}}%
  \vspace{4pt}\par
}

% Subsection styling — timing label
\newcommand{\timing}[1]{%
  \noindent{\small\bfseries\color{pennred}#1}\quad\par\vspace{2pt}%
}

% Slide box environment
\newmdenv[
  backgroundcolor=lightgray,
  linecolor=slidegray,
  linewidth=1pt,
  leftmargin=0pt,
  rightmargin=0pt,
  innerleftmargin=10pt,
  innerrightmargin=10pt,
  innertopmargin=6pt,
  innerbottommargin=6pt,
  skipabove=4pt,
  skipbelow=8pt
]{slidebox}

\setlength{\parindent}{0pt}
\setlength{\parskip}{6pt}

\begin{document}

% =========================================================
% Title block
% =========================================================
\begin{center}
  {\LARGE\bfseries\color{pennblue} Speaking Script}\\[4pt]
  {\large Physics-Grounded Data Augmentation Outperforms Model Adaptation\\
  for Cross-Dataset HAR Transfer}\\[8pt]
  {\normalsize Gengzhan \quad Forest \quad Gengzhan\\
  University of Pennsylvania \quad April 2026}\\[4pt]
  \textcolor{pennred}{\rule{\textwidth}{1.5pt}}\\[2pt]
  {\small Total duration: \textbf{7 minutes} \quad Pace: $\sim$140 words/min}
\end{center}

\vspace{6pt}

% =========================================================
\slidetitle{Slide 1 \ — \ Title}
% =========================================================
\timing{[0:00 – 0:15]}

Hi everyone. We're going to talk about a practical problem in human activity
recognition --- what happens when you train a model on one dataset and try to
deploy it on another. Spoiler: it usually fails badly, and we wanted to understand why.

% =========================================================
\slidetitle{Slide 2 \ — \ Motivation}
% =========================================================
\timing{[0:15 – 1:00]}

So here's the concrete problem. We took UCI~HAR as the source --- that's
the standard waist-mounted sensor dataset --- and tried to run it on WISDM,
which is collected from phones in people's front pockets.

The result: macro F1 of 0.085. That's basically random chance for a 5-class problem.

And this isn't a model quality issue. The same CNN that achieves 95\% on
UCI~HAR completely collapses on WISDM. So the question we asked was:
is this domain shift, and if so, what kind?

% =========================================================
\slidetitle{Slide 3 \ — \ Two Datasets}
% =========================================================
\timing{[1:00 – 1:30]}

These two datasets look similar on the surface --- same activities, same
accelerometer, roughly similar size. But when you look closely, they're actually
quite different.

UCI~HAR is waist-mounted, 50~Hz, and Butterworth low-pass filtered before
release. WISDM is front-pocket, 20~Hz, completely raw.

After physical alignment --- resampling to 20~Hz, converting units, matching
window size --- both datasets give us 51$\times$3 windows.
But the distributions are still very different.

% =========================================================
\slidetitle{Slide 4 \ — \ Seven Sources of Domain Shift}
% =========================================================
\timing{[1:30 – 2:00]}

We identified seven sources of domain shift. I'll go through the top ones.

The dominant source --- and it's not even close --- is gravity axis projection.
Wasserstein distance of \textbf{8.7~m/s$^2$} on the X-axis alone.
That number is basically the entire gravity vector.

The other sources --- amplitude differences, noise structure, class prior shift
--- they all matter, but they're secondary.

% =========================================================
\slidetitle{Slide 5 \ — \ Dominant Shift: Gravity Axis Projection}
% =========================================================
\timing{[2:00 – 3:00]}
\begin{slidebox}
\textcolor{pennred}{\textbf{Key slide. Take your time here.}}
\end{slidebox}

Here's why gravity dominates. In UCI~HAR, the sensor is strapped to the
waist pointing forward. Gravity falls almost entirely along the X-axis ---
mean around $+8.58$~m/s$^2$.

In WISDM, the phone is in your front pocket, free to rotate. Gravity ends up
on the Y-axis, with a lower mean of about $+6.51$~m/s$^2$.

So when a model trains on UCI~HAR, it learns: ``X has a big DC component ---
that's the gravity signal, that's how I tell walking from standing.'' Then on
WISDM, X is near zero-mean. The gravity cue the model learned is just gone.

This isn't a subtle distribution shift. The sensor is physically rotated
roughly 90 degrees.

% =========================================================
\slidetitle{Slide 6 \ — \ Feature Space Gap}
% =========================================================
\timing{[3:00 – 3:30]}

This shows up dramatically in feature space. Before any adaptation, the
multi-kernel MMD between UCI~HAR and WISDM is \textbf{0.830}.

Looking at the PCA projection --- the two datasets are almost completely
disjoint. A model trained on the red cluster genuinely cannot generalize to
the green cluster.

% =========================================================
\slidetitle{Slide 7 \ — \ PhysicalDistortion Pipeline}
% =========================================================
\timing{[3:30 – 4:30]}
\begin{slidebox}
\textcolor{pennred}{\textbf{Key slide. Walk through each operator clearly.}}
\end{slidebox}

So our approach was: instead of adapting the model, correct the data.

We built a five-operator pipeline called \textbf{PhysicalDistortion} that
transforms UCI~HAR training windows to look like they came from a
front-pocket sensor.

\textbf{First}, we apply a 90-degree rotation around the Z-axis, plus some
random wobble to simulate loose pocket placement.

\textbf{Second}, we scale down the gravity component --- 7.29 over 9.69,
which is the empirically measured ratio.

\textbf{Third} --- and this is important --- we apply per-activity, per-axis
amplitude scaling. Not a global scale factor. The ratio for Sitting Y-axis is
4.82$\times$, for Walking Y-axis it's 2.06$\times$. A global scale factor
would get this very wrong.

\textbf{Fourth}, we boost the 0.8 to 2~Hz gait frequency band for locomotion
activities. WISDM has about 40\% more energy there because it's unfiltered.

\textbf{Fifth}, we inject temporally correlated noise --- AR(1) with
coefficient 0.9865 --- because UCI~HAR noise is white after filtering, but
WISDM noise has strong temporal correlation.

Each operator targets a specific, measured physical difference.
No target labels needed.

% =========================================================
\slidetitle{Slide 8 \ — \ After PhysicalDistortion: Distributions Align}
% =========================================================
\timing{[4:30 – 5:00]}

After PhysicalDistortion, the MMD drops from 0.830 to 0.453 ---
a \textbf{45\% reduction}.

You can see in the PCA plot that the augmented UCI data now partially overlaps
with WISDM. The gap is substantially smaller.

Not gone --- that remaining gap is probably why some per-class results are
still imperfect.

% =========================================================
\slidetitle{Slide 9 \ — \ Ablation Results}
% =========================================================
\timing{[5:00 – 5:30]}

Here's the main ablation. Four conditions across five architectures.

C0~Raw is basically non-functional --- average F1 of 0.085. CNN1D with
PhysicalDistortion alone reaches \textbf{0.723}. That's a 6$\times$ improvement
for CNN1D, and 6.8$\times$ on average across all five architectures.

The interesting finding: adding TTBN or DANN on top of PhysicalDistortion
\emph{actually hurts}. Every single model gets worse when you add model-level
adaptation on top of good data augmentation.

% =========================================================
\slidetitle{Slide 10 \ — \ Per-class Analysis: Honest Look}
% =========================================================
\timing{[5:30 – 6:00]}

Looking at per-class performance, the picture is uneven. Sitting at 0.966,
Standing at 0.820, Walking at 0.872 --- these are all solid.

But Upstairs is \textbf{0.445}. That's genuinely bad. Stair climbing produces
a complicated signal with both periodic and aperiodic components, and the
inter-subject variability in how people climb stairs is large.
Window-level augmentation can't model step-to-step variability.
We don't have a good solution for this.

% =========================================================
\slidetitle{Slide 11 \ — \ Why TTBN and DANN Don't Help}
% =========================================================
\timing{[6:00 – 6:20]}

Very briefly --- after PhysicalDistortion, TTBN re-estimates batch norm
statistics from the test batch, but WISDM is 56\% Walking versus UCI's 20\%.
The re-estimated statistics get skewed by the class imbalance.

DANN tries to force domain-invariant features, but the shift is tied to the
activity itself --- gravity projection changes per activity. Forcing invariance
destroys the axis-specific information the classifier actually needs.

% =========================================================
\slidetitle{Slide 12 \ — \ Lessons Learned}
% =========================================================
\timing{[6:20 – 6:45]}

Three takeaways.

\textbf{One}: diagnose before adapting. We spent time measuring the shift
before building anything, and that directly told us what to fix.

\textbf{Two}: physical corrections beat model corrections here. The root cause
was in the data, so fixing the data worked better than fixing the model.

\textbf{Three}: model adaptation made things worse when added on top of good
augmentation. The combination is not always additive.

% =========================================================
\slidetitle{Slide 13 \ — \ Conclusion}
% =========================================================
\timing{[6:45 – 7:00]}

To wrap up: the dominant source of cross-dataset failure between UCI~HAR and
WISDM is gravity axis projection --- a purely physical difference in sensor
placement.

PhysicalDistortion corrects for this without any target labels, taking CNN1D
from F1~0.120 to \textbf{0.723}. The lesson is simple: fix the physics first.
Then think about the model.

Thanks --- happy to take questions.

\vspace{12pt}
\textcolor{pennred}{\rule{\textwidth}{1.5pt}}
\begin{center}
\small\color{gray}
\textit{End of script. Total: $\approx$7:00 at 140 words/min.}
\end{center}

\end{document}
